% Options for packages loaded elsewhere
% Options for packages loaded elsewhere
\PassOptionsToPackage{unicode}{hyperref}
\PassOptionsToPackage{hyphens}{url}
\PassOptionsToPackage{dvipsnames,svgnames,x11names}{xcolor}
%
\documentclass[
  letterpaper,
  DIV=11,
  numbers=noendperiod]{scrartcl}
\usepackage{xcolor}
\usepackage[top=30mm,left=30mm]{geometry}
\usepackage{amsmath,amssymb}
\setcounter{secnumdepth}{-\maxdimen} % remove section numbering
\usepackage{iftex}
\ifPDFTeX
  \usepackage[T1]{fontenc}
  \usepackage[utf8]{inputenc}
  \usepackage{textcomp} % provide euro and other symbols
\else % if luatex or xetex
  \usepackage{unicode-math} % this also loads fontspec
  \defaultfontfeatures{Scale=MatchLowercase}
  \defaultfontfeatures[\rmfamily]{Ligatures=TeX,Scale=1}
\fi
\usepackage{lmodern}
\ifPDFTeX\else
  % xetex/luatex font selection
\fi
% Use upquote if available, for straight quotes in verbatim environments
\IfFileExists{upquote.sty}{\usepackage{upquote}}{}
\IfFileExists{microtype.sty}{% use microtype if available
  \usepackage[]{microtype}
  \UseMicrotypeSet[protrusion]{basicmath} % disable protrusion for tt fonts
}{}
\makeatletter
\@ifundefined{KOMAClassName}{% if non-KOMA class
  \IfFileExists{parskip.sty}{%
    \usepackage{parskip}
  }{% else
    \setlength{\parindent}{0pt}
    \setlength{\parskip}{6pt plus 2pt minus 1pt}}
}{% if KOMA class
  \KOMAoptions{parskip=half}}
\makeatother
% Make \paragraph and \subparagraph free-standing
\makeatletter
\ifx\paragraph\undefined\else
  \let\oldparagraph\paragraph
  \renewcommand{\paragraph}{
    \@ifstar
      \xxxParagraphStar
      \xxxParagraphNoStar
  }
  \newcommand{\xxxParagraphStar}[1]{\oldparagraph*{#1}\mbox{}}
  \newcommand{\xxxParagraphNoStar}[1]{\oldparagraph{#1}\mbox{}}
\fi
\ifx\subparagraph\undefined\else
  \let\oldsubparagraph\subparagraph
  \renewcommand{\subparagraph}{
    \@ifstar
      \xxxSubParagraphStar
      \xxxSubParagraphNoStar
  }
  \newcommand{\xxxSubParagraphStar}[1]{\oldsubparagraph*{#1}\mbox{}}
  \newcommand{\xxxSubParagraphNoStar}[1]{\oldsubparagraph{#1}\mbox{}}
\fi
\makeatother


\usepackage{longtable,booktabs,array}
\newcounter{none} % for unnumbered tables
\usepackage{calc} % for calculating minipage widths
% Correct order of tables after \paragraph or \subparagraph
\usepackage{etoolbox}
\makeatletter
\patchcmd\longtable{\par}{\if@noskipsec\mbox{}\fi\par}{}{}
\makeatother
% Allow footnotes in longtable head/foot
\IfFileExists{footnotehyper.sty}{\usepackage{footnotehyper}}{\usepackage{footnote}}
\makesavenoteenv{longtable}
\usepackage{graphicx}
\makeatletter
\newsavebox\pandoc@box
\newcommand*\pandocbounded[1]{% scales image to fit in text height/width
  \sbox\pandoc@box{#1}%
  \Gscale@div\@tempa{\textheight}{\dimexpr\ht\pandoc@box+\dp\pandoc@box\relax}%
  \Gscale@div\@tempb{\linewidth}{\wd\pandoc@box}%
  \ifdim\@tempb\p@<\@tempa\p@\let\@tempa\@tempb\fi% select the smaller of both
  \ifdim\@tempa\p@<\p@\scalebox{\@tempa}{\usebox\pandoc@box}%
  \else\usebox{\pandoc@box}%
  \fi%
}
% Set default figure placement to htbp
\def\fps@figure{htbp}
\makeatother


% definitions for citeproc citations
\NewDocumentCommand\citeproctext{}{}
\NewDocumentCommand\citeproc{mm}{%
  \begingroup\def\citeproctext{#2}\cite{#1}\endgroup}
\makeatletter
 % allow citations to break across lines
 \let\@cite@ofmt\@firstofone
 % avoid brackets around text for \cite:
 \def\@biblabel#1{}
 \def\@cite#1#2{{#1\if@tempswa , #2\fi}}
\makeatother
\newlength{\cslhangindent}
\setlength{\cslhangindent}{1.5em}
\newlength{\csllabelwidth}
\setlength{\csllabelwidth}{3em}
\newenvironment{CSLReferences}[2] % #1 hanging-indent, #2 entry-spacing
 {\begin{list}{}{%
  \setlength{\itemindent}{0pt}
  \setlength{\leftmargin}{0pt}
  \setlength{\parsep}{0pt}
  % turn on hanging indent if param 1 is 1
  \ifodd #1
   \setlength{\leftmargin}{\cslhangindent}
   \setlength{\itemindent}{-1\cslhangindent}
  \fi
  % set entry spacing
  \setlength{\itemsep}{#2\baselineskip}}}
 {\end{list}}
\usepackage{calc}
\newcommand{\CSLBlock}[1]{\hfill\break\parbox[t]{\linewidth}{\strut\ignorespaces#1\strut}}
\newcommand{\CSLLeftMargin}[1]{\parbox[t]{\csllabelwidth}{\strut#1\strut}}
\newcommand{\CSLRightInline}[1]{\parbox[t]{\linewidth - \csllabelwidth}{\strut#1\strut}}
\newcommand{\CSLIndent}[1]{\hspace{\cslhangindent}#1}



\setlength{\emergencystretch}{3em} % prevent overfull lines

\providecommand{\tightlist}{%
  \setlength{\itemsep}{0pt}\setlength{\parskip}{0pt}}



 


\KOMAoption{captions}{tableheading}
\makeatletter
\@ifpackageloaded{tcolorbox}{}{\usepackage[skins,breakable]{tcolorbox}}
\@ifpackageloaded{fontawesome5}{}{\usepackage{fontawesome5}}
\definecolor{quarto-callout-color}{HTML}{909090}
\definecolor{quarto-callout-note-color}{HTML}{0758E5}
\definecolor{quarto-callout-important-color}{HTML}{CC1914}
\definecolor{quarto-callout-warning-color}{HTML}{EB9113}
\definecolor{quarto-callout-tip-color}{HTML}{00A047}
\definecolor{quarto-callout-caution-color}{HTML}{FC5300}
\definecolor{quarto-callout-color-frame}{HTML}{acacac}
\definecolor{quarto-callout-note-color-frame}{HTML}{4582ec}
\definecolor{quarto-callout-important-color-frame}{HTML}{d9534f}
\definecolor{quarto-callout-warning-color-frame}{HTML}{f0ad4e}
\definecolor{quarto-callout-tip-color-frame}{HTML}{02b875}
\definecolor{quarto-callout-caution-color-frame}{HTML}{fd7e14}
\makeatother
\makeatletter
\@ifpackageloaded{caption}{}{\usepackage{caption}}
\AtBeginDocument{%
\ifdefined\contentsname
  \renewcommand*\contentsname{Table of contents}
\else
  \newcommand\contentsname{Table of contents}
\fi
\ifdefined\listfigurename
  \renewcommand*\listfigurename{List of Figures}
\else
  \newcommand\listfigurename{List of Figures}
\fi
\ifdefined\listtablename
  \renewcommand*\listtablename{List of Tables}
\else
  \newcommand\listtablename{List of Tables}
\fi
\ifdefined\figurename
  \renewcommand*\figurename{Figure}
\else
  \newcommand\figurename{Figure}
\fi
\ifdefined\tablename
  \renewcommand*\tablename{Table}
\else
  \newcommand\tablename{Table}
\fi
}
\@ifpackageloaded{float}{}{\usepackage{float}}
\floatstyle{ruled}
\@ifundefined{c@chapter}{\newfloat{codelisting}{h}{lop}}{\newfloat{codelisting}{h}{lop}[chapter]}
\floatname{codelisting}{Listing}
\newcommand*\listoflistings{\listof{codelisting}{List of Listings}}
\makeatother
\makeatletter
\makeatother
\makeatletter
\@ifpackageloaded{caption}{}{\usepackage{caption}}
\@ifpackageloaded{subcaption}{}{\usepackage{subcaption}}
\makeatother
\usepackage{bookmark}
\IfFileExists{xurl.sty}{\usepackage{xurl}}{} % add URL line breaks if available
\urlstyle{same}
% fallback for those not using the hyperref driver hyperxmp:
\makeatletter
\@ifundefined{xmpquote}{\newcommand{\xmpquote}[1]{#1}}{}
\makeatother
\hypersetup{
  pdftitle={Preregistration part II: benefits{,} barriers and building up skills},
  pdfauthor={Sara Lil Middleton},
  colorlinks=true,
  linkcolor={blue},
  filecolor={Maroon},
  citecolor={Blue},
  urlcolor={Blue},
  pdfcreator={LaTeX via pandoc}}


\title{Preregistration part II: benefits, barriers and building up
skills}
\author{Sara Lil Middleton}
\date{12/11/2026}
\begin{document}
\maketitle


\subsection{Prerequisites}\label{prerequisites}

\begin{tcolorbox}[enhanced jigsaw, arc=.35mm, bottomrule=.15mm, bottomtitle=1mm, breakable, colback=white, colbacktitle=quarto-callout-important-color!10!white, colframe=quarto-callout-important-color-frame, coltitle=black, left=2mm, leftrule=.75mm, opacityback=0, opacitybacktitle=0.6, rightrule=.15mm, title=\textcolor{quarto-callout-important-color}{\faExclamation}\hspace{0.5em}{Important}, titlerule=0mm, toprule=.15mm, toptitle=1mm]

Before completing this session, please carefully read about the
prerequisites.

\end{tcolorbox}

\begin{itemize}
\item
  Completion of preregistration training materials (preregistration part
  I {[}coming soon!{]}
\item
  Access to Build up your preregistration toolkit! {[}coming soon{]}
\end{itemize}

\textbf{Presenter Notes}: These are the \textbf{presenter notes}. You
will find a script for the presenter for every slide. In presentation
mode, your audience will not be able to see these presenter notes, they
are only visible to the presenter.

For this session, it is important that you are familiar with what a
preregistration is, since this session is building up on that.

\textbf{Instructor Notes}: There are also \textbf{instructor notes}. For
some slides, there will be pedagogical tips, suggestions for activities
and troubleshooting tips for issues your audience might run into. You
can find these notes underneath the presenter notes.

These are the prerequisites for this session. Before you get started on
this session with your audience, you need to ensure that the audience
fulfills these criteria.

\begin{center}\rule{0.5\linewidth}{0.5pt}\end{center}

\subsection{Your study preregistration
journey}\label{your-study-preregistration-journey}

From: \emph{``I am not sure what some of the opportunities are or how to
overcome some of the challenges of doing a preregistration} To:
``\emph{``I feel able discuss the benefits and barriers of
preregistration and identify which part of a preregistration I can
improve on''}

\textbf{Presenter Notes}: In part 2 of this preregistration workshop I
hope you will go from not being sure of the opportunities and challenges
of doing a preregistration to being able to discuss benefits and
barriers and identify an area of the preregistration process you can
improve on.

\begin{center}\rule{0.5\linewidth}{0.5pt}\end{center}

\subsection{Recap of key concepts in preregistration part
I}\label{recap-of-key-concepts-in-preregistration-part-i}

Study preregistration is an open science practice that:

\begin{itemize}
\item
  \textbf{safe guards} researchers from \textbf{questionable research
  practices} and \textbf{biases} by making decision-making processes
  more explicit.
\item
  involves uploading a time-stamped, publicly-available* \textbf{record}
  of a research study plan onto a register, which is given a unique and
  permanent study identifier.
\item
  *Preregistrations do not need to be made public immediately and can be
  embargoed for a period of years
  (e.g.~\href{https://help.osf.io/article/634-addressing-fear-of-scooping-on-osf-preregistrations}{Open
  Science Framework preregistrations} can be embargoed for 4 years from
  first submission)
\end{itemize}

\textbf{Presenter Notes}: Study preregistration is an open science
practice that safe guard researchers from questionable research
practices and biases by making decision-making processes more explicit.

\begin{itemize}
\item
  It involves uploading a time-stamped, publicly-available record of a
  research study plan onto a register, which is given a unique and
  permanent study identifier.
\item
  Preregistrations do not need to be made public immediately and can be
  embargoed for a period of years. For example OSF preregistration can
  be embargoed for a period of up to 4 years.
\end{itemize}

\textbf{Instructor Notes}: A quick recap from preregistration part 1
helps learners situate themselves in the preregistration learning
journey and before a short questions session (next slide).

\begin{center}\rule{0.5\linewidth}{0.5pt}\end{center}

\subsection{Questions from preregistration part
I}\label{questions-from-preregistration-part-i}

\textbf{Presenter Notes}: Before we get started with preregistration
part II, now is a good time to clarify questions from part I.

\textbf{Instructor Notes}: This slide is dedicated to clarifying
questions from the previous submodule and/or to discuss assignments. -
Additional slides may need to be added depending on the nature of the
homework assignments. - Critical for the learning process to ensure that
learners are on the same page and have been able to achieve the learning
goals of the previous workshop. - Not applicable if this set of slides
corresponds to the first submodule of a new module.

\begin{center}\rule{0.5\linewidth}{0.5pt}\end{center}

\subsection{Before we start: Survey
time!}\label{before-we-start-survey-time}

\subsubsection{Question 1}

\textbf{What is your level of familiarity with the key components of a
preregistration (e.g., elements found in a preregistration document )?}

\begin{itemize}
\item[$\square$]
  1. I don't know
\item[$\square$]
  2. I know a few key components
\item[$\square$]
  3. I have a good grasp of the key components
\item[$\square$]
  4. I am very familiar/have done a preregistration before
\end{itemize}

\subsubsection{Question 2}

\textbf{How would you rate your confidence to carry out a
preregistration for your study?} (1 to 5 scale: \textbf{1} = Not
confident at all, \textbf{2} = Slightly confident, \textbf{3} = Somewhat
confident, \textbf{4} = Very confident, \textbf{5} = Completely
confident)

\begin{itemize}
\item[$\square$]
  1. Not confident at all
\item[$\square$]
  2. Slightly confident
\item[$\square$]
  3. Somewhat confident
\item[$\square$]
  4. Very confident
\item[$\square$]
  5. Completely confident
\end{itemize}

\begin{center}\rule{0.5\linewidth}{0.5pt}\end{center}

\subsection{Discussion of survey
results}\label{discussion-of-survey-results}

What do we see in the results?

\textbf{Presenter Notes}: Let us have a look at these results.

\textbf{Instructor Notes}:\textbf{Aim}: Briefly examine the answers
given to each question interactively with the group. Use visuals from
the survey to highlight specific answers.

\textbf{Accessibility Tip}: Have people work in pairs if someone did not
bring their phone/does not have a phone. This survey also cannot be
completed in an asynchronous setting.

\begin{center}\rule{0.5\linewidth}{0.5pt}\end{center}

\subsection{Learning goals}\label{learning-goals}

\textbf{Aim}: Examine the benefits of preregistration to researchers and
wider research ecosystem, and discuss common challenges and how to
address them.

There are \textbf{two} sections to this submodule:

\begin{itemize}
\item
  \textbf{Section 1}: Benefits and barriers of preregistration
\item
  \textbf{Section 2:} Building up skills in preregistration
\end{itemize}

\textbf{Presenter Notes}: The overall aim of this submodule is to
examine the benefits of preregistration to researchers and wider
research ecosystem, and discuss common challenges and how to address
them.

There are 2 sections: first we will discuss the benefits and barriers of
doing a preregistration and then in the second section we will look at
practicing improving our preregistration skills.

\textbf{Instructor Notes}: Giving the whole overview at the start of the
session provides learners with structure and clarity about what to
expect.

\begin{center}\rule{0.5\linewidth}{0.5pt}\end{center}

\subsection{Section 1: Learning goals and
overview}\label{section-1-learning-goals-and-overview}

Section 1 is all about recognizing common benefits and barriers of
preregistration for researchers and wider research ecosystem.

Activities: Quiz and self-reflection exercise

After completing section 1 you should be able to:

\begin{itemize}
\tightlist
\item
  \textbf{Explain} the perceived barriers to doing a preregistration
\item
  \textbf{Identify} key benefits of preregistration
\end{itemize}

\textbf{Presenter Notes}: We will begin by outlining the common benefits
and barriers of preregistration for researchers and wider research
ecosystem.

There will be a quiz at the end and along the way several
self-reflection exercises.

After completing section 1 you should be able to:

\begin{itemize}
\tightlist
\item
  \textbf{Explain} the perceived barriers to doing a preregistration
\item
  \textbf{Identify} key benefits of preregistration
\end{itemize}

\begin{center}\rule{0.5\linewidth}{0.5pt}\end{center}

\subsection{Section 2: Learning goals and
overview}\label{section-2-learning-goals-and-overview}

Section 2 is about practicing improving a component of a study
preregistration.

Activity: Build up your preregistration toolkit! (individual work)

After completing section 2 you should be able to:

\begin{itemize}
\tightlist
\item
  \textbf{Self-reflect} on challenging parts of the preregistration
  process
\item
  \textbf{Practice} improving elements of a preregistration
\end{itemize}

\textbf{Presenter Notes}: In section 2 we turn our attention to building
up our registration toolkit, by practicing improving a component of a
study preregistration.

\begin{center}\rule{0.5\linewidth}{0.5pt}\end{center}

\subsection{Section 1: Benefits and barriers of
preregistration}\label{section-1-benefits-and-barriers-of-preregistration}

\textbf{Presenter Notes}: In section 1, we will go over the key benefits
and barriers of preregistration. We will first consider how benefits and
opportunities of preregistration exist from individual researchers all
the way to the wider research ecosystem.

\begin{center}\rule{0.5\linewidth}{0.5pt}\end{center}

\subsection{Benefits for researchers
I}\label{benefits-for-researchers-i}

Improved planning and project management

\begin{itemize}
\tightlist
\item
  The process of filling out preregistration document helps
  \textbf{identify key information} needed at study design stage
\item
  Outlining data collection methods and analysis plans which
  \textbf{reduces ambiguity} and \textbf{errors} during data analysis
  stage
\end{itemize}

Builds your reputation

\begin{itemize}
\tightlist
\item
  Transparently show to the wider research community your plans and
  tests
\item
  Signal that you have taken steps to \textbf{mitigate} against
  hindsight and confirmation \textbf{bias} (Wagenmakers \& Dutilh,
  2016).
\end{itemize}

\textbf{Presenter Notes}: A key area that preregistration benefits
researchers is by improving the planning and project management as time
is spent outlining data collection methods and analysis plans which
reduces ambiguity and errors during data analysis stage. The process of
filling out preregistration document helps identify key information
needed at study design stage, which is particularly helpful for early
career researchers.

Preregistrations also allow you to build your reputation by -
Transparently showing to the wider research community your plans and
tests and - signalling that you have taken steps to mitigate against
hindsight and confirmation bias.

\begin{center}\rule{0.5\linewidth}{0.5pt}\end{center}

\subsection{Benefits for researchers
II}\label{benefits-for-researchers-ii}

Easier to write up and publish findings

\begin{itemize}
\item
  Front loading the thinking and planning makes \textbf{writing up}
  results and discussion a smoother process
\item
  If you preregister your analytic \textbf{code} on \textbf{simulated
  data} the results of your confirmatory analyses will be
  \textbf{generated quicker} once you have collected and cleaned data
\item
  If preregistration is approved as a registered report during stage 1
  journal peer review, you have \textbf{``in-principle acceptance''}
\end{itemize}

Visible productivity

\begin{itemize}
\tightlist
\item
  Publicly sharing outputs that others can \textbf{evaluate} beyond a
  journal article (as evaluation of outputs can be influenced by journal
  prestige)
\end{itemize}

\textbf{Presenter Notes}: Preregistration makes it easier to write up
and publish findings by front loading the thinking and planning. This
makes writing up results and discussion a smoother process!

If you preregister your analytic \textbf{code} on \textbf{simulated
data} the results of your confirmatory analyses will be
\textbf{generated quicker} once you have collected and cleaned data

And, if preregistration is approved as a registered report during stage
1 journal peer review, you have ``in-principle acceptance''

Closely linked, is the benefit of more visibility of your work beyond
journal articles (as evaluation of outputs can be influenced by journal
prestige).

\begin{center}\rule{0.5\linewidth}{0.5pt}\end{center}

\subsection{Benefits for teams}\label{benefits-for-teams}

Clearer involvement of collaborators/project partners

\begin{itemize}
\tightlist
\item
  Provides opportunities for project partners to discuss and
  \textbf{refine plans} before data collection
\item
  Detailed research plans provide \textbf{clarity} for project partners,
  which is increasingly important in the era of \textbf{large-scale
  collaborations}
\end{itemize}

\textbf{Presenter Notes}: Clearer involvement of collaborators/project
partners is another important benefit of preregistration at the level of
teams.

Preregistration provides opportunities for project partners to discuss
and refine plans before data collection - a chance to catch any
potential errors early!

Detailed research plans provide clarity for project partners, which is
increasingly important in the era of large-scale collaborations (also
known as Big Team Science).

\begin{center}\rule{0.5\linewidth}{0.5pt}\end{center}

\subsection{Benefits to wider research
ecosystem}\label{benefits-to-wider-research-ecosystem}

Increase research credibility and bias

\begin{itemize}
\tightlist
\item
  \textbf{Mitigates} against questionable research practices
\item
  Offers more complete reporting and \textbf{transparency} of research
  project
\item
  \textbf{Reduces publication bias}, as preregistration/registered
  report is available to read (if not embargoed)
\end{itemize}

Check thinking

\begin{itemize}
\tightlist
\item
  Other people reading your findings can \textbf{follow your thought
  process}
\item
  Gain a deeper understanding of a project and make an \textbf{informed
  assessment} (Binney et al., 2025)
\end{itemize}

\textbf{Presenter Notes}: Thinking at the scale of academic disciplines
and wider academia, preregistration helps in two main ways:

firstly by increasing research credibility by - Mitigating against
questionable research practices and - Offering more complete reporting
and transparency of research project. - by mitigating against the
publication bias or file drawer effect as the project plans are
available and accessible on a repository or published as a registered
report even if there are no statistically significant results. This is
especially valuable when it comes to meta-analyses.

Secondly, it allows other people to check thinking by - Allowing other
people to read your findings and follow your thought process enabling
them to - gain a deeper understanding of a project.

\textbf{Instructor Notes}: Note about research credibility in slide at
the end of section 1 about PARKing (``Preregistrating after the results
are known'').

\begin{center}\rule{0.5\linewidth}{0.5pt}\end{center}

\subsection{Benefits recap}\label{benefits-recap}

\begin{itemize}
\item
  There are many benefits to doing a preregistration for individual
  researchers, teams and the wider research ecosystem.
\item
  Overall, by taking the time to carefully plan your study design, data
  collection and approach to analysis, research becomes more rigorous,
  transparent and credible.
\end{itemize}

\begin{tcolorbox}[enhanced jigsaw, arc=.35mm, bottomrule=.15mm, bottomtitle=1mm, breakable, colback=white, colbacktitle=quarto-callout-note-color!10!white, colframe=quarto-callout-note-color-frame, coltitle=black, left=2mm, leftrule=.75mm, opacityback=0, opacitybacktitle=0.6, rightrule=.15mm, title=\textcolor{quarto-callout-note-color}{\faInfo}\hspace{0.5em}{Reflection point}, titlerule=0mm, toprule=.15mm, toptitle=1mm]

Do you relate to the benefits of preregistration mentioned? Are there
any other benefits you can think of?

\end{tcolorbox}

\begin{tcolorbox}[enhanced jigsaw, arc=.35mm, bottomrule=.15mm, bottomtitle=1mm, breakable, colback=white, colbacktitle=quarto-callout-tip-color!10!white, colframe=quarto-callout-tip-color-frame, coltitle=black, left=2mm, leftrule=.75mm, opacityback=0, opacitybacktitle=0.6, rightrule=.15mm, title=\textcolor{quarto-callout-tip-color}{\faLightbulb}\hspace{0.5em}{Find out more}, titlerule=0mm, toprule=.15mm, toptitle=1mm]

More benefits about preregistration and registered reports from a
student, author and journal editor perspective are outlined by Gya et
al. (2023).

\end{tcolorbox}

\textbf{Presenter Notes}: Before we move on the considering some of the
barriers of preregistration, The take home message is that by taking the
time to carefully plan your study design, data collection and approach
to analysis research becomes more rigorous, transparent and credible.

You can explore this topic more, you can read additional benefits of
preregistration from multiple perspectives in the field of ecology and
evolution by Gya and colleagues.

\textbf{Instructor Notes}: Reflection point can be a time for learners
to quietly consolidate the points (1-2 minutes) covered on benefits of
preregistration and consider their own contexts. This reflection
exercise can be expanded so learners can share their reflections in
class.

\begin{center}\rule{0.5\linewidth}{0.5pt}\end{center}

\subsection{Barriers to
preregistration}\label{barriers-to-preregistration}

Wider context:

\begin{itemize}
\item
  There are significant challenges present in the current
  \textbf{research} and \textbf{assessment} \textbf{culture}
  (e.g.~``publish or perish'') which influences the uptake of
  preregistration.
\item
  Researchers have \textbf{different knowledge} and \textbf{experience}
  levels with Open Science practices that reflect in \textbf{attitudes}
  and \textbf{adoption} of preregistration.
\end{itemize}

\textbf{Presenter Notes}: In the final part of section 1 we will outline
some common challenges researchers new to the practice of
preregistration can encounter.

Some context to keep in mind:

\begin{itemize}
\item
  There are significant challenges present in the current research and
  assessment culture (e.g.~``publish or perish'') which influences the
  uptake of preregistration.
\item
  Researchers have different knowledge and experience levels with Open
  Science practices that reflect in attitudes and adoption of
  preregistration.
\end{itemize}

\begin{center}\rule{0.5\linewidth}{0.5pt}\end{center}

\subsection{Challenges: time}\label{challenges-time}

\textbf{Time}: \emph{``I haven't got time to do a preregistration''}

\begin{itemize}
\tightlist
\item
  \emph{Reframing:} front-loaded effort during the planning stage can
  help later on during write up and during peer review (remember the
  benefits to researchers slide earlier)
\end{itemize}

\begin{tcolorbox}[enhanced jigsaw, arc=.35mm, bottomrule=.15mm, bottomtitle=1mm, breakable, colback=white, colbacktitle=quarto-callout-note-color!10!white, colframe=quarto-callout-note-color-frame, coltitle=black, left=2mm, leftrule=.75mm, opacityback=0, opacitybacktitle=0.6, rightrule=.15mm, title=\textcolor{quarto-callout-note-color}{\faInfo}\hspace{0.5em}{Did you know?}, titlerule=0mm, toprule=.15mm, toptitle=1mm]

\href{https://rr.peercommunityin.org/PCIRegisteredReports/}{Peer
Community in Registered Reports} is an initiative where the research
community reviews and recommends registered reports. There is the option
to speed up the Stage 1 peer review of a registered report with the
`scheduled review'
\href{https://rr.peercommunityin.org/help/guide_for_reviewers\#h_61998243643551613309672490}{track}.
Authors initially submit a one-page ``snapshot'' instead of a full stage
1 report.

\end{tcolorbox}

\textbf{Presenter Notes}: A common argument against the uptake of study
preregistration is

``I haven't got time to do a study registration'' to which a
counterargument would be: front-loaded effort during the planning stage
can help later on during write up and during peer review.

To make you aware there is an initiative called ``Peer Community in
Registered Reports'', where the research community reviews and
recommends registered reports. There is the option to speed up the Stage
1 peer review of a registered report with the `scheduled review' track
Authors initially submit a one-page ``snapshot'' instead of a full stage
1 report.

\textbf{Instructor Notes}: It is important for instructors to
acknowledge the wider `publish or perish' research culture. There is a
some truth to the statement, as study planning should take time. It can
be helpful to reframe preregistrations as an upfront time investment
that will pay off later (link also to benefits of preregistration
e.g.~easier to write up findings).

\begin{center}\rule{0.5\linewidth}{0.5pt}\end{center}

\subsection{Challenges: statistical
confidence}\label{challenges-statistical-confidence}

\textbf{Statistical confidence}: \emph{``How do I know my planned
statistical analysis is the right one?''}

\begin{itemize}
\tightlist
\item
  \emph{Strategy:} Simulate random data with the distribution and sample
  size of expected real data and run the statistical test planned for
  the study.
\item
  Repeat the statistical test (e.g.~100 x) and check results are within
  the false-positive finding rate of 5\%
\item
  Run a power analysis by simulating data with the expected effect size
  for the study and check model registers effect at least 80\% of the
  time
\end{itemize}

\begin{tcolorbox}[enhanced jigsaw, arc=.35mm, bottomrule=.15mm, bottomtitle=1mm, breakable, colback=white, colbacktitle=quarto-callout-note-color!10!white, colframe=quarto-callout-note-color-frame, coltitle=black, left=2mm, leftrule=.75mm, opacityback=0, opacitybacktitle=0.6, rightrule=.15mm, title=\textcolor{quarto-callout-note-color}{\faInfo}\hspace{0.5em}{Additional resource}, titlerule=0mm, toprule=.15mm, toptitle=1mm]

\href{https://lmu-osc.github.io/teaching-material-for-novices/open-research-tools/simulations/}{Training
materials} for novices: Introduction to simulations and power analysis
in R.

\end{tcolorbox}

\textbf{Presenter Notes}: Another common barrier is lack of confidence
in choosing the right statistical test. The strategy would be to run
simulations and power analysis. Simulate random data with the
distribution and sample size of expected real data and run the
statistical test planned for the study. - Repeat the statistical test
(e.g.~100 x) and check results are within the false-positive finding
rate of 5\% - Run a power analysis by simulating data with the expected
effect size for the study and check model registers effect at least 80\%
of the time

There are introductory guides on this topic which you can look at in
your own time.

\begin{center}\rule{0.5\linewidth}{0.5pt}\end{center}

\subsection{Challenges: inflexible}\label{challenges-inflexible}

\textbf{Lack of flexibility:} \emph{``Preregistration keeps you in a
box''}

\begin{itemize}
\tightlist
\item
  \emph{Reframing:} Preregistrations are a ``plan not a prison'' as you
  can document any deviations from your plan when you write up the
  research findings (DeHaven, 2017).
\item
  Preregistrations can be tailored to the study type (e.g incremental
\item
  transparently documenting exploratory-confirmatory pathway or
  sequential - multi-stage contingent experiments).
\end{itemize}

\begin{tcolorbox}[enhanced jigsaw, arc=.35mm, bottomrule=.15mm, bottomtitle=1mm, breakable, colback=white, colbacktitle=quarto-callout-note-color!10!white, colframe=quarto-callout-note-color-frame, coltitle=black, left=2mm, leftrule=.75mm, opacityback=0, opacitybacktitle=0.6, rightrule=.15mm, title=\textcolor{quarto-callout-note-color}{\faInfo}\hspace{0.5em}{Find out more}, titlerule=0mm, toprule=.15mm, toptitle=1mm]

Binney et al. (2025) discuss practical guidance of balancing mitigating
biases with flexibility in the research process. Lakens (2024) discusses
when and how to deviate from a preregistration.

\end{tcolorbox}

\textbf{Presenter Notes}: ``Preregistration keeps you in a box'' is
another common perceived barrier. There is different degrees of
flexibility in preregistations as you can document any deviations from
your plan when you write up the research findings. It is important to
describe changes to the design, sampling or analysis plans clearly.

Preregistration can be tailored to the study type. They can be
incremental (e.g.~transparently documenting exploratory-confirmatory
pathway) or sequential (e.g.~multi-stage contingent experiments).

\begin{center}\rule{0.5\linewidth}{0.5pt}\end{center}

\subsection{Challenges: scooping?}\label{challenges-scooping}

\textbf{Fear of scooping}: \emph{``What if someone steals my ideas
before they are published?''}

\begin{itemize}
\item
  \emph{Counterargument:} Preregistrations do not need to be made public
  right away and can be embargoed for a period of years
  (e.g.~\href{https://help.osf.io/article/634-addressing-fear-of-scooping-on-osf-preregistrations}{Open
  Science Framework preregistrations} can be embargoed for 4 years from
  first submission)
\item
  The embargo period allows researchers time after preregistering their
  study to conduct the research, analyse, write up and publish their
  findings.
\end{itemize}

\textbf{Presenter Notes}: A commonly cited barrier to preregistration is
tied to intellectual property with the fear of being scooped. The
pressure to publish contributes to the fear of having research ideas
``stolen'' before they are published, as preregistrations are publically
available.

To counteract this barrier, preregistered studies can be uploaded onto
preregistration platforms and made private for a period of time. The
Open Science Framework allows an embargo period of 4 years from first
submission before the preregistration is made public.

The embargo period allows researchers time after preregistering their
study to conduct the research, analyse, write up and publish their
findings.

\begin{center}\rule{0.5\linewidth}{0.5pt}\end{center}

\subsection{Challenges: incentives}\label{challenges-incentives}

\textbf{Lack of incentives}: \emph{``I don't see the point''}

\begin{itemize}
\tightlist
\item
  \emph{Counterargument:} aside from individual researcher benefits
  (increased reflection, better planning, reduced likelihood of
  systematic errors), journals are increasingly prioritizing
  preregistrations.
\end{itemize}

\begin{tcolorbox}[enhanced jigsaw, arc=.35mm, bottomrule=.15mm, bottomtitle=1mm, breakable, colback=white, colbacktitle=quarto-callout-note-color!10!white, colframe=quarto-callout-note-color-frame, coltitle=black, left=2mm, leftrule=.75mm, opacityback=0, opacitybacktitle=0.6, rightrule=.15mm, title=\textcolor{quarto-callout-note-color}{\faInfo}\hspace{0.5em}{Find out more}, titlerule=0mm, toprule=.15mm, toptitle=1mm]

Next time you read a journal article in your field, take a look at the
journal's webpage and see if they accept preregistered studies and if
there are guidelines for publishing registered reports. This information
is usually in the author resources/guidelines section.

\end{tcolorbox}

\textbf{Presenter Notes}: The counterargument to ``I don't see the
point'' that aside from individual researcher benefits (increased
reflection, better planning, reduced likelihood of systematic errors),
journals are increasingly prioritizing study registrations.

A nice little exercise to do next time you read a journal article is to
see if they accept preregistered studies and if there are guidelines for
publishing registered reports. This information is usually in the author
resources/guidelines section.

\textbf{Instructor Notes}: This additional exercise is optional and can
be set as a homework. It is designed for learners to see that prevalence
and acceptance of preregistration is growing. Psychology journals are
likely have higher prevalance see: Simonsohn (2023)
https://datacolada.org/115.

\begin{center}\rule{0.5\linewidth}{0.5pt}\end{center}

\subsection{Challenges: overwhelm}\label{challenges-overwhelm}

\textbf{Overwhelm}: \emph{``There is so much that goes into a
preregistration''}

\begin{itemize}
\tightlist
\item
  \emph{Counterargument:} templates and FAQ pages can guide you and if
  you are unable to officially register a study, adopting principles of
  study preregistration is a good step (e.g.~Open Science Buffet
  approach (Bergmann, 2023).
\end{itemize}

\begin{tcolorbox}[enhanced jigsaw, arc=.35mm, bottomrule=.15mm, bottomtitle=1mm, breakable, colback=white, colbacktitle=quarto-callout-note-color!10!white, colframe=quarto-callout-note-color-frame, coltitle=black, left=2mm, leftrule=.75mm, opacityback=0, opacitybacktitle=0.6, rightrule=.15mm, title=\textcolor{quarto-callout-note-color}{\faInfo}\hspace{0.5em}{Find out more}, titlerule=0mm, toprule=.15mm, toptitle=1mm]

The Open Science Framework (OSF) has a
\href{https://help.osf.io/article/330-welcome-to-registrations\#FAQs-GG_d_}{frequently
asked questions} page on preregistrations and registered reports.

\end{tcolorbox}

\textbf{Presenter Notes}: When something is new it can be hard to wrap
your head around the concept or practice and it can feel overwhelming.
With preregistration, there are lots of templates that guide you through
the process and support pages.

We will revisit how to address this challenge in section 2!

\begin{center}\rule{0.5\linewidth}{0.5pt}\end{center}

\subsection{Challenges: lack of
support}\label{challenges-lack-of-support}

\textbf{Lack of support}: \emph{``No one in my team knows or cares about
preregistration''}

\begin{itemize}
\item
  \emph{Strategy:} Often a reluctance to try out preregistration comes
  from a place of lack of familiarity with the process.
\item
  Sharing resources, and suggest a group meeting to discuss/learn more
  about preregistration to increase familiarity.
\item
  Reframe preregistration as a creative opportunity to plan great
  science! Planning is where the science is, as other research
  processes, like data collection can be delegated
\end{itemize}

\begin{tcolorbox}[enhanced jigsaw, arc=.35mm, bottomrule=.15mm, bottomtitle=1mm, breakable, colback=white, colbacktitle=quarto-callout-note-color!10!white, colframe=quarto-callout-note-color-frame, coltitle=black, left=2mm, leftrule=.75mm, opacityback=0, opacitybacktitle=0.6, rightrule=.15mm, title=\textcolor{quarto-callout-note-color}{\faInfo}\hspace{0.5em}{Additional resource}, titlerule=0mm, toprule=.15mm, toptitle=1mm]

Relate to this challenge? See the preregistration
\href{https://zenodo.org/records/22307709}{case study} which provides
strategies.

\end{tcolorbox}

\textbf{Presenter Notes}: A challenge which can be a bit more complex to
resolve, is when there is a lack of support from supervisors or group
members. This can come from a place of lack of familiarity. It is
important to understand where this resistance is coming from. Opening up
a discussion with the wider group can be beneficial place to start.

\textbf{Instructor Notes}: You can expand on student-supervisor
reframing of preregistration: planning is where the science really
happens, a lot of other tasks can be delegated (data collection,
statistical analyses, writing up), but not the thinking /creative and
smart planning. Students can make a plan and then have their supervisors
check the robustness of their plans.

Additional resource with case study can be set as an in class exercise
or homework (time allocation 15-60 minutes, depending on how in-depth
you wish to go).

\begin{center}\rule{0.5\linewidth}{0.5pt}\end{center}

\subsection{What preregistration is
not}\label{what-preregistration-is-not}

Preregistration does not prevent fraud or research misconduct (e.g.
PARKing), but it does make it harder.

\begin{tcolorbox}[enhanced jigsaw, arc=.35mm, bottomrule=.15mm, bottomtitle=1mm, breakable, colback=white, colbacktitle=quarto-callout-note-color!10!white, colframe=quarto-callout-note-color-frame, coltitle=black, left=2mm, leftrule=.75mm, opacityback=0, opacitybacktitle=0.6, rightrule=.15mm, title=\textcolor{quarto-callout-note-color}{\faInfo}\hspace{0.5em}{Definition}, titlerule=0mm, toprule=.15mm, toptitle=1mm]

\textbf{PARKing}: ``Preregistrating after the results are known'' is
considered research misconduct where researchers complete an experiment
before preregistering which invalidates the purpose of preregistration
(Parsons et al., 2022).

\end{tcolorbox}

\textbf{Presenter Notes}: It is important to clarify what
preregistrations don't do.

They do not prevent fraud or research misconduct (e.g.~PARKing).

PARKing or ``Preregistrating after the results are known'' is considered
research misconduct where researchers complete an experiment before
preregistering which invalidates the purpose of preregistration.

\textbf{Instructor Notes}: It is important to communicate with learners
that preregistration is not the ``silver bullet'' that can resolve.

\begin{center}\rule{0.5\linewidth}{0.5pt}\end{center}

\subsection{Challenges recap}\label{challenges-recap}

\begin{itemize}
\item
  The lack of time and incentives, feelings of overwhelm and lack of
  support are common barriers that often overlap when in it comes to
  preregistration.
\item
  There is lots of support available and you are not alone!
\end{itemize}

\begin{tcolorbox}[enhanced jigsaw, arc=.35mm, bottomrule=.15mm, bottomtitle=1mm, breakable, colback=white, colbacktitle=quarto-callout-tip-color!10!white, colframe=quarto-callout-tip-color-frame, coltitle=black, left=2mm, leftrule=.75mm, opacityback=0, opacitybacktitle=0.6, rightrule=.15mm, title=\textcolor{quarto-callout-tip-color}{\faLightbulb}\hspace{0.5em}{Find out more}, titlerule=0mm, toprule=.15mm, toptitle=1mm]

Have a browse of some of the challenges and strategies to address
barriers to preregistration here {[}Preregistration tutorial coming
soon!{]}

\end{tcolorbox}

\textbf{Presenter Notes}: Let us recap the challenges we have discussed.
- The lack of time and incentives, feelings of overwhelm and lack of
support are common barriers that often overlap when in it comes to
preregistration.

\begin{itemize}
\tightlist
\item
  There is lots of support available and you are not alone!
\end{itemize}

\begin{center}\rule{0.5\linewidth}{0.5pt}\end{center}

\subsection{Quiz}\label{quiz}

\textbf{Presenter Notes}: To round off section 1 we have a short quiz of
3 questions where the answers will appear after each question.

\begin{center}\rule{0.5\linewidth}{0.5pt}\end{center}

\textbf{Q1. Which options below are benefits of registered reports}

\begin{itemize}
\tightlist
\item[$\square$]
  In-principle acceptance (IPA) of manuscripts
\item[$\square$]
  Prevents fraud
\item[$\square$]
  Mitigates publication bias
\item[$\square$]
  Opportunity for community feedback on project plans
\end{itemize}

\textbf{Presenter Notes}: Out of the following four options, which ones
are benefits of doing a registered report?

\begin{itemize}
\tightlist
\item
  In-principle acceptance (IPA) of manuscripts
\item
  Prevents fraud
\item
  Mitigates publication bias
\item
  Opportunity for community feedback on project plans
\end{itemize}

\begin{center}\rule{0.5\linewidth}{0.5pt}\end{center}

\textbf{A1. Which option(s) below are benefits of registered reports}

\begin{itemize}
\tightlist
\item[$\boxtimes$]
  In-principle acceptance (IPA) of manuscripts
\item[$\square$]
  Prevents fraud
\item[$\boxtimes$]
  Mitigates publication bias
\item[$\boxtimes$]
  Opportunity for community feedback on project plans
\end{itemize}

\begin{tcolorbox}[enhanced jigsaw, arc=.35mm, bottomrule=.15mm, bottomtitle=1mm, breakable, colback=white, colbacktitle=quarto-callout-note-color!10!white, colframe=quarto-callout-note-color-frame, coltitle=black, left=2mm, leftrule=.75mm, opacityback=0, opacitybacktitle=0.6, rightrule=.15mm, title=\textcolor{quarto-callout-note-color}{\faInfo}\hspace{0.5em}{Reflecting on answer}, titlerule=0mm, toprule=.15mm, toptitle=1mm]

Although registered reports (and preregistrations) mitigate against
questionable research practices, they do not prevent research misconduct
or fraud (Yamada, 2018).

\end{tcolorbox}

\textbf{Presenter Notes}: The answers are: - In-principle acceptance
(IPA) of manuscripts - Mitigates publication bias - Opportunity for
community feedback on project plans

Although registered reports (and preregistrations) mitigate against
questionable research practices, they do not prevent research misconduct
or fraud.

\begin{center}\rule{0.5\linewidth}{0.5pt}\end{center}

\textbf{Q2. Which from the options below is an example of how
preregistrations benefit the wider research ecosystem level?}

\begin{itemize}
\tightlist
\item[$\square$]
  Mitigates publication bias
\item[$\square$]
  Improves planning and project management
\item[$\square$]
  Builds up reputation
\item[$\square$]
  Visible productivity
\end{itemize}

\textbf{Presenter Notes}: Out of the following four options, which ones
are benefits of doing a registered report?

\begin{itemize}
\tightlist
\item
  Mitigates publication bias
\item
  Improves planning and project management
\item
  Builds up reputation
\item
  Visible productivity
\end{itemize}

\begin{center}\rule{0.5\linewidth}{0.5pt}\end{center}

\textbf{A2. Which from the options below is an example of how
preregistrations benefit the wider research ecosystem level?}

\begin{itemize}
\tightlist
\item[$\boxtimes$]
  Mitigates publication bias
\item[$\square$]
  Improves planning and project management
\item[$\square$]
  Builds up reputation
\item[$\square$]
  Visible productivity
\end{itemize}

\begin{tcolorbox}[enhanced jigsaw, arc=.35mm, bottomrule=.15mm, bottomtitle=1mm, breakable, colback=white, colbacktitle=quarto-callout-note-color!10!white, colframe=quarto-callout-note-color-frame, coltitle=black, left=2mm, leftrule=.75mm, opacityback=0, opacitybacktitle=0.6, rightrule=.15mm, title=\textcolor{quarto-callout-note-color}{\faInfo}\hspace{0.5em}{Reflecting on answer}, titlerule=0mm, toprule=.15mm, toptitle=1mm]

Improves planning and project management, building up reputation and
visible productivity are examples of benefits at the individual
researcher level.

\end{tcolorbox}

\textbf{Presenter Notes}: The answer is mitigating against publication,
which is an example of a benefit of preregistration at the wider
research ecosystem level, particularly when it comes to meta-research.

\begin{center}\rule{0.5\linewidth}{0.5pt}\end{center}

\textbf{Q3. What can I do if I feel overwhelmed or lack time to do a
preregistration?}

\begin{itemize}
\tightlist
\item[$\square$]
  Give up on trying out a preregistration
\item[$\square$]
  Find a preregistration template
\item[$\square$]
  Reading Frequently Asked Question pages
\item[$\square$]
  Trying an Open Science Buffet approach (Bergmann, 2023)
\end{itemize}

\textbf{Presenter Notes}: What can I do if I feel overwhelmed or lack
time to do a preregistration?**

\begin{itemize}
\tightlist
\item[$\square$]
  Give up on trying out a preregistration
\item[$\square$]
  Find a preregistration template
\item[$\square$]
  Reading Frequently Asked Question pages
\item[$\square$]
  Trying an Open Science Buffet approach
\end{itemize}

\begin{center}\rule{0.5\linewidth}{0.5pt}\end{center}

\textbf{A3. What can I do if I feel overwhelmed or lack time to do a
preregistration?}

\begin{itemize}
\tightlist
\item[$\square$]
  Give up on trying out a preregistration
\item[$\boxtimes$]
  Find a preregistration template
\item[$\boxtimes$]
  Reading Frequently Asked Question pages
\item[$\boxtimes$]
  Trying an Open Science Buffet approach (Bergmann, 2023)
\end{itemize}

\begin{tcolorbox}[enhanced jigsaw, arc=.35mm, bottomrule=.15mm, bottomtitle=1mm, breakable, colback=white, colbacktitle=quarto-callout-note-color!10!white, colframe=quarto-callout-note-color-frame, coltitle=black, left=2mm, leftrule=.75mm, opacityback=0, opacitybacktitle=0.6, rightrule=.15mm, title=\textcolor{quarto-callout-note-color}{\faInfo}\hspace{0.5em}{Reflecting on answer}, titlerule=0mm, toprule=.15mm, toptitle=1mm]

When trying out a new way of working, it is common for things to feel
overwhelming. The Open Science Framework (OSF) has a
\href{https://help.osf.io/article/330-welcome-to-registrations\#FAQs-GG_d_}{frequently
asked questions} page on preregistrations and registered reports.

\end{tcolorbox}

\textbf{Presenter Notes}: The answer is the bottom three answers: - Find
a preregistration template - Reading Frequently Asked Question pages -
Trying an Open Science Buffet approach

\begin{center}\rule{0.5\linewidth}{0.5pt}\end{center}

\subsection{Summary of section 1}\label{summary-of-section-1}

\begin{itemize}
\tightlist
\item
  Taking the time to \textbf{carefully plan} your study design, data
  collection and approach to analysis, leads to more \textbf{rigorous},
  \textbf{transparent} and \textbf{credible research}.\\
\item
  Lack of time or incentives are two common barriers to adopting the
  practice of study preregistration
\item
  It is important to openly discuss and address challenges
\item
  Preregistration does not prevent fraud or research misconduct (e.g.
  PARKing).
\end{itemize}

\begin{tcolorbox}[enhanced jigsaw, arc=.35mm, bottomrule=.15mm, bottomtitle=1mm, breakable, colback=white, colbacktitle=quarto-callout-note-color!10!white, colframe=quarto-callout-note-color-frame, coltitle=black, left=2mm, leftrule=.75mm, opacityback=0, opacitybacktitle=0.6, rightrule=.15mm, title=\textcolor{quarto-callout-note-color}{\faInfo}\hspace{0.5em}{Reflection point}, titlerule=0mm, toprule=.15mm, toptitle=1mm]

Which benefit and barrier resonated with you most?

\end{tcolorbox}

\textbf{Presenter Notes}: To summarize section 1:

\begin{itemize}
\item
  Preregistration leads to more rigorous, transparent and credible
  research by taking the time to carefully plan the study design, data
  collection and approach to analysis.
\item
  Lack of time or incentives are two common barriers to adopting the
  practice of study registration
\item
  Important to discuss and address challenges
\item
  Preregistration does not prevent fraud or research misconduct (e.g.
  PARKing).
\end{itemize}

\begin{center}\rule{0.5\linewidth}{0.5pt}\end{center}

\subsection{Break! 15 minutes}\label{break-15-minutes}

\begin{center}\rule{0.5\linewidth}{0.5pt}\end{center}

\subsection{Section 2: Build up your preregistration
toolkit!}\label{section-2-build-up-your-preregistration-toolkit}

\textbf{Presenter Notes}: Now that you've taken a look at some of the
challenges and opportunities of preregistration let's practice improving
elements of preregistration.

In the spirit of an ``Open Science buffet'' (Bergmann, 2023), let's
start to build up your study registration toolkit!

Starting small can help manage the feeling of overwhelm with trying out
a practice for the first time!

Remember that one of the common challenges is feelings of overwhelm with
doing a preregistration, and uncertainty about where to start.

\begin{center}\rule{0.5\linewidth}{0.5pt}\end{center}

\subsection{Main preregistration
elements}\label{main-preregistration-elements}

Recap of the main elements that make up a preregistration document:

\begin{itemize}
\tightlist
\item[$\square$]
  Research questions and hypotheses
\item[$\square$]
  Sampling plan
\item[$\square$]
  Design plan
\item[$\square$]
  Outliers and exceptions
\item[$\square$]
  Analysis plan
\end{itemize}

\begin{tcolorbox}[enhanced jigsaw, arc=.35mm, bottomrule=.15mm, bottomtitle=1mm, breakable, colback=white, colbacktitle=quarto-callout-tip-color!10!white, colframe=quarto-callout-tip-color-frame, coltitle=black, left=2mm, leftrule=.75mm, opacityback=0, opacitybacktitle=0.6, rightrule=.15mm, title=\textcolor{quarto-callout-tip-color}{\faLightbulb}\hspace{0.5em}{Tip}, titlerule=0mm, toprule=.15mm, toptitle=1mm]

If you need a reminder of the essential and recommended elements that
make up a preregistration document see preregistration part I (LINK).

\end{tcolorbox}

\textbf{Presenter Notes}:

In the previous preregistration workshop, we outlined the parts that go
into a preregistration document. As a reminder they are:

\begin{itemize}
\tightlist
\item
  Research questions and hypotheses
\item
  Sampling plan
\item
  Design plan
\item
  Outliers and exceptions
\item
  Analysis plan
\end{itemize}

\begin{center}\rule{0.5\linewidth}{0.5pt}\end{center}

\subsection{Exercise instructions}\label{exercise-instructions}

The objective of the build up your preregistration toolkit exercise is
to:

\begin{itemize}
\tightlist
\item
  \textbf{Identify} what makes preregistration elements vague
\item
  \textbf{Practice} improving the precision and clarity of
  preregistration elements
\end{itemize}

Steps:

\begin{enumerate}
\def\labelenumi{\arabic{enumi}.}
\tightlist
\item
  \emph{Group practice} on one example preregistration snippet
\item
  \emph{Work individually} on worksheet with 8 snippets (40 minutes)
\end{enumerate}

\emph{Presenter Notes}:

The two main objectives of the build up your preregistration toolkit
exercise is to:

\begin{itemize}
\tightlist
\item
  Identify what makes preregistration elements vague
\item
  Practice improving the precision and clarity of preregistration
  elements
\end{itemize}

We will first do a practice round as a group then we will work on our
own

\textbf{Instructor Notes}: Remind learners that even if they have
familiarity and practice in writing these preregistration elements, it
can be helpful to revisit examples outside your own research area.

Doing the group exercise first provides scaffolding for learners to
practice together and see a modeled answer before doing the exercise on
their own. It also helps set up expectations and an opportunity to
clarify understanding.

\begin{center}\rule{0.5\linewidth}{0.5pt}\end{center}

\subsection{Group exercise}\label{group-exercise}

A research team wants to do a study on the topic of how indoor plants
impact a person's mood. Take a look at this research question:

\textbf{How do indoor plants benefit people?}

\begin{enumerate}
\def\labelenumi{\arabic{enumi}.}
\tightlist
\item
  On your own, \emph{think} about how the question can be improved
\item
  \emph{Discuss} answer with neighbor and \emph{write} out a joint
  answer
\item
  We will \emph{share} our answers in class
\end{enumerate}

\begin{tcolorbox}[enhanced jigsaw, arc=.35mm, bottomrule=.15mm, bottomtitle=1mm, breakable, colback=white, colbacktitle=quarto-callout-tip-color!10!white, colframe=quarto-callout-tip-color-frame, coltitle=black, left=2mm, leftrule=.75mm, opacityback=0, opacitybacktitle=0.6, rightrule=.15mm, title=\textcolor{quarto-callout-tip-color}{\faLightbulb}\hspace{0.5em}{Guiding questions}, titlerule=0mm, toprule=.15mm, toptitle=1mm]

What is ambiguous about this question? How would you measure
``benefit''? How would you study the interaction between plants and
people?

\end{tcolorbox}

\textbf{Presenter Notes:} Let's have a go at improving this answer: A
research team wants to do a study on the topic of how indoor plants
impact a person's mood.

First think about it on your own, then discuss answer with neighbor.
Then after a couple of minutes we will share some answers as a whole
group.

Take a look at the guiding questions to help you.

\begin{center}\rule{0.5\linewidth}{0.5pt}\end{center}

\subsection{Group exercise: answer}\label{group-exercise-answer}

\emph{From:} How do indoor plants benefit people?

\emph{To:} How does having physical contact with indoor plants reduce
psychological and physiological stress in young-adult office workers?

\begin{tcolorbox}[enhanced jigsaw, arc=.35mm, bottomrule=.15mm, bottomtitle=1mm, breakable, colback=white, colbacktitle=quarto-callout-note-color!10!white, colframe=quarto-callout-note-color-frame, coltitle=black, left=2mm, leftrule=.75mm, opacityback=0, opacitybacktitle=0.6, rightrule=.15mm, title=\textcolor{quarto-callout-note-color}{\faInfo}\hspace{0.5em}{Note}, titlerule=0mm, toprule=.15mm, toptitle=1mm]

Notice how the improved answer states the type of interaction between
plants and people, more specific details about types of benefits of
interest and the target demographic (e.g.~young-adult office workers).

\end{tcolorbox}

\textbf{Presenter Notes}: So we went from ``how do indoor plants benefit
people?'' to a clearer research question of:

How does having physical contact with indoor plants reduce psychological
and physiological stress in young-adult office workers?

Notice how the improved answer states the type of interaction between
plants and people, more specific details about types of benefits of
interest and the target demographic (e.g.~young-adult office workers).

\begin{center}\rule{0.5\linewidth}{0.5pt}\end{center}

\subsection{Now it is your turn!}\label{now-it-is-your-turn}

\begin{enumerate}
\def\labelenumi{\arabic{enumi}.}
\tightlist
\item
  \textbf{Access} worksheet {[}LINK{]}
\item
  \textbf{Read} instructions
\item
  \textbf{Work} through improving the snippets (40 minutes)
\end{enumerate}

\textbf{Presenter Notes}: Now it is your turn!

Access the ``build up your preregistration toolkit! worksheet here and
follow the instructions. We will have a bout 40 minutes to do this
before we wrap up.

\textbf{Instructor Notes}: Check learners have access to the worksheet.
Pause to take questions.

\begin{center}\rule{0.5\linewidth}{0.5pt}\end{center}

\subsection{Exercise reflections}\label{exercise-reflections}

\textbf{Presenter Notes}: We have some time now to discuss how the
exercise went for you, which relates to the assignment.

\textbf{Instructor Notes}: Invite learners to share the experience of
doing this exercise. This slide acts a bridge to the assignment (next
slide)

\begin{center}\rule{0.5\linewidth}{0.5pt}\end{center}

\subsection{Assignment}\label{assignment}

One step at a time, then tell a buddy!

\begin{enumerate}
\def\labelenumi{\arabic{enumi}.}
\tightlist
\item
  \emph{Self-reflect} and commit to taking one action related to the
  preregistration process
\item
  Once you have chosen an action, \emph{tell a friend} for
  accountability!
\end{enumerate}

\begin{tcolorbox}[enhanced jigsaw, arc=.35mm, bottomrule=.15mm, bottomtitle=1mm, breakable, colback=white, colbacktitle=quarto-callout-tip-color!10!white, colframe=quarto-callout-tip-color-frame, coltitle=black, left=2mm, leftrule=.75mm, opacityback=0, opacitybacktitle=0.6, rightrule=.15mm, title=\textcolor{quarto-callout-tip-color}{\faLightbulb}\hspace{0.5em}{Possible actions}, titlerule=0mm, toprule=.15mm, toptitle=1mm]

Writing out clearer hypotheses, researching a study registration
template for your subject area. See ``additional resources'' in build up
your preregistration toolkit worksheet.

\end{tcolorbox}

\textbf{Presenter Notes}: From the main exercise the worksheet you may
have identified a particular part of the preregistration process you
wish to improve.

In this extension exercise you are invited to self-reflect and commit to
taking one action related to the preregistration process (e.g.~writing
out clearer hypotheses, researching a study registration template for
your subject area).

Once you have chosen an action, tell a friend for accountability! To
help you, take a look at the ``additional resources'' section.

\textbf{Instructor Notes}: This extension exercise is designed to be an
empowering and manageable action to commit to, which links to the spirit
of ``Open Science buffet'' and the journey outlined throughout this
session. It also promotes self-reflection, critical thinking and a sense
of agency in the choices that are made in how research is carried out.

\begin{center}\rule{0.5\linewidth}{0.5pt}\end{center}

\subsection{To conclude: Survey time!}\label{to-conclude-survey-time}

\textbf{Presenter Notes}: Before we wrap up with a final summary, let's
revisit the survey questions from the beginning.

\textbf{Instructor Notes}:

\begin{itemize}
\item
  \textbf{Aim}: This post-submodule survey serves to examine learners'
  current knowledge about the sumodule's topic.
\item
  Use free survey software such as particify or formR to establish the
  following questions (shown on separate slides)
\item
  Use identical questions as in the pre-submodule survey to be able to
  directly compare
\end{itemize}

\begin{center}\rule{0.5\linewidth}{0.5pt}\end{center}

\textbf{1. What is your level of familiarity with the key components of
a preregistration (e.g., elements found in a preregistration document
)?}

\begin{itemize}
\item[$\square$]
  1. I don't know
\item[$\square$]
  2. I know a few key components
\item[$\square$]
  3. I have a good grasp of the key components
\item[$\square$]
  4. I am very familiar/have done a preregistration before
\end{itemize}

\begin{center}\rule{0.5\linewidth}{0.5pt}\end{center}

\textbf{2. How would you rate your confidence to carry out a
preregistration for your study?} (1 to 5 scale: \textbf{1} = Not
confident at all, \textbf{2} = Slightly confident, \textbf{3} = Somewhat
confident, \textbf{4} = Very confident, \textbf{5} = Completely
confident)

\begin{itemize}
\item[$\square$]
  1. Not confident at all
\item[$\square$]
  2. Slightly confident
\item[$\square$]
  3. Somewhat confident
\item[$\square$]
  4. Very confident
\item[$\square$]
  5. Completely confident
\end{itemize}

\begin{center}\rule{0.5\linewidth}{0.5pt}\end{center}

\subsection{Final Summary}\label{final-summary}

\begin{itemize}
\item
  Preregistration and registered reports offer many \textbf{benefits} to
  researchers, teams and the wider community (e.g.~\textbf{building
  reputation} and more \textbf{rigorous}, \textbf{transparent} and
  \textbf{credible research}).
\item
  A number of \textbf{challenges} exist to the wider uptake of
  preregistration (e.g.~perceived inflexibility, time commitment,
  overwhelm).
\item
  Adopting a \textbf{few elements} of study preregistration at a time is
  good approach to working your way up completing a preregistration or
  registered report.
\item
  There are lots of places to get \textbf{support} (e.g.~Open Science
  Framework community pages).
\end{itemize}

\textbf{Presenter Notes}: Preregistration and registered reports offer
many benefits to researchers, teams and the wider community
(e.g.~building reputation and more rigorous, transparent and credible
research).

\begin{itemize}
\item
  A number of challenges exist to the wider uptake of preregistration
  (e.g.~perceived inflexibility, time commitment, overwhelm).
\item
  Adopting a few elements of study preregistration at a time is good
  approach to working your way up completing a preregistration or
  registered report.
\item
  There are lots of places to get support (e.g.~Open Science Framework
  community pages).
\end{itemize}

\textbf{Instructor Notes}: Allow time for questions and remind learners
of their assignment.

\begin{center}\rule{0.5\linewidth}{0.5pt}\end{center}

\subsection{Additional resources}\label{additional-resources}

\href{https://www.youtube.com/watch?v=77YeUe_Bkco}{Discussion} about
merits vs limits of study preregistration from IEEE Visualization
Conference 2022

\subsection{References}\label{references}

\protect\phantomsection\label{refs}
\begin{CSLReferences}{1}{0}
\bibitem[\citeproctext]{ref-bergmannBuffetApproachOpen2023}
Bergmann, C. (2023). \emph{The {Buffet Approach} to {Open Science}}.

\bibitem[\citeproctext]{ref-binneyPracticalRoutesPreregistration2025}
Binney, R. J., Smith, L. J., Rossit, S., Demeyere, N., Learmonth, G.,
Olgiati, E., Halai, A. D., Rounis, E., Evans, J., Edelstyn, N. M. J., \&
McIntosh, R. D. (2025). Practical routes to preregistration: A guide to
enhanced transparency and rigour in neuropsychological research.
\emph{Brain Communications}, \emph{7}(3), fcaf162.
\url{https://doi.org/10.1093/braincomms/fcaf162}

\bibitem[\citeproctext]{ref-dehavenPreregistrationPlanNot2017}
DeHaven, A. C. (2017). Preregistration: {A Plan}, {Not} a {Prison}. In
\emph{Preregistration: A Plan, Not a Prison}.

\bibitem[\citeproctext]{ref-gyaRegisteredReportsNew2023}
Gya, R., Birkeli, K., Dahle, I. J., Foote, C. G., Geange, S. R., Lynn,
J. S., Töpper, J. P., Vandvik, V., Zernichow, C., \& Jenkins, G. B.
(2023). Registered {Reports}: {A} new chapter at {\emph{Ecology}}{
\emph{\&} }{\emph{Evolution}}. \emph{Ecology and Evolution},
\emph{13}(4), e10023. \url{https://doi.org/10.1002/ece3.10023}

\bibitem[\citeproctext]{ref-lakensWhenHowDeviate2024}
Lakens, D. (2024). When and {How} to {Deviate From} a {Preregistration}.
\emph{Collabra: Psychology}, \emph{10}(1), 117094.
\url{https://doi.org/10.1525/collabra.117094}

\bibitem[\citeproctext]{ref-parsonsCommunitysourcedGlossaryOpen2022}
Parsons, S., Azevedo, F., Elsherif, M. M., Guay, S., Shahim, O. N.,
Govaart, G. H., Norris, E., O'Mahony, A., Parker, A. J., Todorovic, A.,
Pennington, C. R., Garcia-Pelegrin, E., Lazić, A., Robertson, O.,
Middleton, S. L., Valentini, B., McCuaig, J., Baker, B. J., Collins, E.,
\ldots{} Aczel, B. (2022). A community-sourced glossary of open
scholarship terms. \emph{Nature Human Behaviour}, \emph{6}(3), 312--318.
\url{https://doi.org/10.1038/s41562-021-01269-4}

\bibitem[\citeproctext]{ref-wagenmakersSevenSelfishReasons2016}
Wagenmakers, E.-J., \& Dutilh, G. (2016). \emph{Seven {Selfish Reasons}
for {Preregistration}}.

\bibitem[\citeproctext]{ref-yamadaHowCrackPreregistration2018}
Yamada, Y. (2018). How to {Crack Pre}-registration: {Toward Transparent}
and {Open Science}. \emph{Frontiers in Psychology}, \emph{9}, 1831.
\url{https://doi.org/10.3389/fpsyg.2018.01831}

\end{CSLReferences}

\begin{center}\rule{0.5\linewidth}{0.5pt}\end{center}

\subsection{Licence and citation
statement}\label{licence-and-citation-statement}

This current work by Sara Lil Middleton and Malika Ihle is licensed
under a
\href{https://creativecommons.org/licenses/by-sa/4.0/deed.en}{Creative
Commons Attribution 4.0 International SA License} (CC-BY-SA 4.0).

Please cite these materials as: Middleton, L. S. \& Ihle, M. (2026).
\emph{Preregistration part II: benefits, barriers and building up
skills} {[}Quarto slides{]}. \href{https://www.osc.lmu.de/}{LMU Open
Science Center} and \href{https://forrt.org/}{FORRT}.
\url{https://lmu-osc.github.io/teaching-material-for-novices/open-research-principles/preregistration-part-II}

\textbf{Presenter Notes}: The
\href{https://creativecommons.org/licenses/by-sa/4.0/deed.en}{Creative
Commons Attribution 4.0 International SA License}, or CC BY-SA 4.0,
allows others to copy, share, and adapt a work in any medium, including
for commercial purposes. These permissions are broad and cannot be
withdrawn as long as the license terms are followed. The main
requirement is attribution: users must give appropriate credit to the
original creator, provide a link to the license, and clearly indicate
whether any changes were made, without implying endorsement by the
original author. In addition, the ShareAlike condition means that if
someone modifies or builds upon the work, the resulting material must be
distributed under the same CC BY-SA 4.0 license, or a compatible one.
Finally, users are not allowed to apply legal or technical restrictions,
such as DRM, that would prevent others from exercising these same
rights.

Code snippets are dedicated to the public domain and licensed under a
CC0 1.0
\href{https://creativecommons.org/publicdomain/zero/1.0/}{Creative
Commons Universal License}. You may use, modify, distribute, and sell
the code snippets for any purpose, without permission or attribution.
The code snippets are provided ``as is'', without warranty of any kind.

\begin{center}\rule{0.5\linewidth}{0.5pt}\end{center}

\subsection{Contribution statement}\label{contribution-statement}

\textbf{Creator}: Middleton, Sara Lil
(\pandocbounded{\includegraphics[keepaspectratio]{images/ORCID-iD_icon_16x16.png}}\href{https://orcid.org/0000-0001-5307-8029}{0000-0001-5307-8029})

\textbf{Reviewer}: Ihle, Malika
(\pandocbounded{\includegraphics[keepaspectratio]{images/ORCID-iD_icon_16x16.png}}
\href{https://orcid.org/0000-0002-3242-5981}{0000-0002-3242-5981})

\textbf{Acknowledgments}: Gracia Prum and Dejana Damjanovic




\end{document}
